% =====================================================================
% QUESTION 1 — abstract classes & interfaces  (4 + 6 = 10 marks)
% Demonstrates: \javafile, \splitcode, \expectedoutput
% =====================================================================
\question{1}{4 + 6 = 10}

\subquestion{a}{4}
\textbf{Find and correct the errors.} The file \texttt{Shape.java} below is
meant to declare an abstract class \texttt{Shape} that implements the
\texttt{Comparable} interface. It contains \textbf{three} errors. For each
error, state what is wrong and write the corrected line.

\begin{javafile}{Shape.java}
public abstract class Shape implements Comparable<Shape> {
    protected String name;

    public Shape(String name) {
        this.name = name;
    }

    // ERROR 1: an abstract method must not have a body
    public abstract double area() {
        return 0.0;
    }

    // ERROR 2: a method declaration is missing its parentheses
    abstract void printName;

    // ERROR 3: an overriding method cannot reduce visibility
    int compareTo(Shape other) {
        return Double.compare(area(), other.area());
    }
}
\end{javafile}

\subquestion{b}{6}
The interface \texttt{Operation} and its implementation \texttt{Calculator}
are given on the left. Complete the code on the right so that the program
prints the result of \texttt{apply(8, 12)}.

\begin{splitcode}{Given Code}{Complete the Code}
\begin{lstlisting}
interface Operation {
    int apply(int a, int b);
}

class Calculator implements Operation {
    @Override
    public int apply(int a, int b) {
        return a + b;
    }
}
\end{lstlisting}
\splitdivider
\begin{lstlisting}
public class Main {
    public static void main(String[] args) {
        // 1. Declare an Operation reference that
        //    points to a new Calculator object.
        // 2. Call apply(8, 12) and print the result.
        //
        // WRITE YOUR CODE HERE
    }
}
\end{lstlisting}
\end{splitcode}

\vspace{0.2cm}
\noindent When completed correctly, the program must produce the console
output shown below.

\begin{expectedoutput}
> Result: 20
\end{expectedoutput}
